% !TEX root = ../report.tex

\section*{Abstract}

This paper examines how dependence across financial assets changes between calm and stressed markets. We study daily returns for 14 exchange-traded funds covering equities, bonds, commodities, gold, oil, and real estate. A two-state Markov-switching model on monthly S\&P~500 returns identifies the primary regimes, validated against a VIX-and-drawdown rule and a weekly specification. Regime-conditional correlations are evaluated using bootstrap inference, principal component analysis, and a Forbes--Rigobon volatility adjustment.

The results are panel-dependent. Average correlation rises strongly among international equities but falls slightly across asset classes, where tighter risk-asset comovement is offset by lower correlations involving government bonds and gold. Nevertheless, both panels become more concentrated in stress, with a stronger dominant component and fewer effective independent directions. After adjustment, no pair retains a robust correlation breakdown, indicating that the raw increases largely reflect higher stress-period volatility. Diversification therefore weakens most clearly within international equities, while sovereign bonds and gold remain comparatively resilient.

\noindent\textit{Keywords:} asset correlation; market stress; Markov switching; bootstrap inference; principal component analysis; Forbes--Rigobon adjustment; diversification

\noindent\textit{Abbreviations:} ETF, exchange-traded fund; PCA, principal component analysis; VIX, volatility index

\newpage
